\documentclass[aps,prd,twocolumn,nofootinbib,superscriptaddress,floatfix]{revtex4-2}

\usepackage{amsmath,amssymb,amsfonts,mathtools}
\usepackage{graphicx}
\usepackage{booktabs}
\usepackage{bm}
\usepackage{microtype}
\usepackage{placeins}
\usepackage[colorlinks=true,citecolor=blue,linkcolor=blue,urlcolor=blue]{hyperref}
\pdfstringdefDisableCommands{\def\\{ }}

\newcommand{\ii}{\mathrm{i}}
\newcommand{\dd}{\mathrm{d}}
\newcommand{\bx}{\bm{x}}
\newcommand{\by}{\bm{y}}
\newcommand{\bA}{\bm{A}}
\newcommand{\bG}{\bm{\Gamma}}
\newcommand{\bH}{\bm{H}}
\newcommand{\cG}{\mathcal{G}}
\newcommand{\cZ}{\mathcal{Z}}
\newcommand{\inner}[2]{\left\langle #1\middle|#2\right\rangle}

\begin{document}

\title{Relative binning for gravitational-wave microlensing\\
in the Chang--Refsdal model}

\author{Tejaswi Venumadhav}
\affiliation{Department of Physics, University of California, Santa Barbara, California 93106, USA}
\date{July 15, 2026}

\begin{abstract}
Stars in a lensing galaxy can microlens a strongly lensed gravitational wave
and imprint diffraction and interference fringes across the detector band.
Relative binning accelerates gravitational-wave parameter estimation by
representing nearby waveforms through smoothly varying ratios that need to be
evaluated at only a sparse set of frequencies.  Applied directly to a
microlensed waveform, this method loses its advantage because small changes in
the lens parameters shift the fringes and make the ratio rapidly varying.  We
reorganize the Chang--Refsdal amplification factor as a coherent sum in which
the image-delay phases are kept explicit and the remaining complex amplitudes
vary slowly with frequency.  Away from caustics these amplitudes reduce to the
individual image contributions.  At moderate frequencies and through folds
and cusps, we use the full wave-optics result to correct the separated
contributions while preserving their sum exactly.  We evaluate the total
amplification from the analytic point-mass solution through source derivatives
and include external convergence by an exact rescaling, without numerical
integration over the lens plane.  In the examples studied here, interpolating
the separated amplitudes while treating the delay phases analytically reduces
the required frequency sampling by more than an order of magnitude compared
with direct interpolation of the full amplification ratio.
\end{abstract}

\maketitle

\section{Introduction}
\label{sec:introduction}

Gravitational waves retain phase coherence over the frequency band of an
interferometric detector.  A gravitational lens can therefore produce
observable diffraction and interference when the delay between competing
propagation paths is comparable to the inverse wave frequency
\cite{NakamuraDeguchi1999,TakahashiNakamura2003}.  Stellar-mass objects lie in
this regime for ground-based detectors.  Their influence is enhanced when the
source is also strongly lensed by a galaxy, since the stellar population in
the lens galaxy perturbs each macro image and can imprint frequency-dependent
modulations on the signal
\cite{ChristianVitaleLoeb2018,MeenaBagla2020,MishraEtAl2021}.  On the scale of
a single microlens Einstein radius, the smooth galaxy potential is described by
a constant convergence and shear.  A point mass in this local field is the
Chang--Refsdal lens \cite{ChangRefsdal1979,ChangRefsdal1984,AnEvans2006}.

Bayesian parameter estimation requires repeated evaluations of the
gravitational-wave likelihood.  Relative binning accelerates this calculation
by choosing a waveform near the high-likelihood region as a fiducial model,
expressing each candidate as a frequency-dependent ratio to that model, and
evaluating the likelihood from precomputed sums on a sparse frequency grid
\cite{ZackayDaiVenumadhav2018,DaiVenumadhavZackay2018,KrishnaEtAl2023,NarolaEtAl2024}.
For nearby compact-binary waveforms, the amplitude and phase of this ratio vary
slowly with frequency, so only a modest number of frequency samples is needed.

A wave-optics amplification factor does not share this property when it is
treated as a single function.  The signal is a coherent sum of contributions
with different image delays.  Changing the lens parameters moves the
interference fringes, and the candidate-to-fiducial ratio can develop rapid
oscillations and sharp features near the zeros of the fiducial amplification.
This occurs even when the candidate and fiducial sources remain within the
same two-image or four-image region.  A change in image topology creates a
second difficulty, but it is not the origin of the failure of ordinary
relative binning.

The stationary-phase expansion suggests how to reorganize the calculation.
Each geometrical image contributes a phase fixed by its time delay, multiplied
by a factor that changes much more slowly with frequency.  We keep the delay
phases separate and interpolate the accompanying factors.  The isolated-image
expansion alone is not accurate at moderate frequency or near a caustic, so we
use it only to define the resolved-image limit.  At finite frequency we assign
the difference between its coherent sum and the full wave-optics result to the
same separated contributions.  Their sum then reproduces the full
amplification exactly, while each contribution approaches the corresponding
image in the high-frequency limit.

A decomposition tied directly to the physical images cannot be continued
through a fold or cusp.  To obtain a representation that is valid throughout
the positive-parity Chang--Refsdal caustic network, we keep four labels at all
parameter values.  Labels without a real image are placed at the nearest
critical point.  When images are created, they inherit these labels
continuously.  An algebraic projection enforces the exact total amplification
for any mixture of real and virtual labels.

The paper develops the representation in the order in which its physical
motivation arises.  Section~\ref{sec:waveoptics} introduces the amplification
factor and the analytic source-plane operator used to evaluate it.
Section~\ref{sec:stationary} relates the observed fringes to interference
between geometrical images.  Section~\ref{sec:rbfailure} shows why direct
relative binning fails.  Section~\ref{sec:localchannels} constructs separated
image contributions within a region of fixed topology.  Folds, cusps, and the
global four-label extension are discussed in
Secs.~\ref{sec:topology} and~\ref{sec:globalchannels}.  The corresponding
multi-component relative-binning likelihood is derived in
Sec.~\ref{sec:multichannel-rb}.  Numerical tests are presented in
Sec.~\ref{sec:numerics}.

\section{Wave optics in the Chang--Refsdal lens}
\label{sec:waveoptics}

\subsection{Fermat potential and amplification factor}

The Fermat potential gives the dimensionless propagation time for a ray that
crosses the lens plane at $\bx$ when the unlensed source position is $\by$.
Its stationary points are the geometrical images.  The wave-optics amplitude is
obtained by coherently summing over the complete lens plane
\cite{SchneiderEhlersFalco1992,NakamuraDeguchi1999}.

We measure angles in the Einstein-radius units of the point mass and write
\begin{equation}
 \tau(\bx;\by)
 =\frac12\bx^{\mathsf T}\bA\bx-\by\mathbin{\cdot}\bx
  +\frac12|\by|^2-\ln|\bx|,
 \label{eq:fermat}
\end{equation}
where
\begin{equation}
 \bA=(1-\kappa)\bm{1}-\bG,
 \label{eq:macro-matrix}
\end{equation}
and
\begin{equation}
 \bG=\gamma
 \begin{pmatrix}
  \cos 2\beta & \sin 2\beta\\
  \sin 2\beta & -\cos 2\beta
 \end{pmatrix}.
 \label{eq:shear-matrix}
\end{equation}
The parameters $\kappa$, $\gamma$, and $\beta$ are the external convergence,
shear magnitude, and shear orientation.  We consider a positive-parity macro
image,
\begin{equation}
 1-\kappa>|\gamma|.
 \label{eq:positive-parity}
\end{equation}
Macro saddle images require a different treatment of the large-scale
oscillatory integral and are outside the present scope.

With the standard Fresnel normalization, the amplification factor is
\begin{equation}
 F(w,\by)=\frac{w}{2\pi\ii}
 \int_{\mathbb R^2}\dd^2x\,
 \exp\!\left[\ii w\tau(\bx;\by)\right].
 \label{eq:diffraction-integral}
\end{equation}
The normalization includes the magnification of the local macro field.  The
dimensionless frequency is
\begin{equation}
 w=\frac{8\pi G M_{\rm L}(1+z_{\rm L})}{c^3}f,
 \label{eq:dimensionless-frequency}
\end{equation}
for a point mass $M_{\rm L}$ at redshift $z_{\rm L}$.  A lensed
frequency-domain waveform is
\begin{equation}
 \widetilde h_{\rm L}(f)=F[w(f),\by]\,\widetilde h_{\rm U}(f).
 \label{eq:lensed-waveform}
\end{equation}
A constant added to $\tau$ changes only the overall arrival time.  We subtract
the minimum image delay in the figures and numerical comparisons.

Figure~\ref{fig:amplification} shows representative amplification curves in
the two-image and four-image regions and near a fold and cusp.  The low-
frequency portion is plotted on a logarithmic frequency axis.  The
high-frequency portion is sampled and displayed linearly because the
geometrical-optics fringes are approximately periodic in $w$.  Sampling those
fringes uniformly in $\log w$ produces visual aliasing.

\begin{figure*}[t]
 \centering
 \includegraphics[width=0.96\textwidth]{figures/amplification_mixed.pdf}
 \caption{Magnitude of the Chang--Refsdal amplification factor for
 representative source positions at $\kappa=\gamma=0.2$ and $\beta=0$.
 The low-frequency segment of each panel uses a logarithmic horizontal axis;
 the high-frequency segment uses a linear axis.  The source positions are $\by=(0.55,0)$, $(0.10,0.10)$,
 $(-0.1389,0.1793)$, and $(-0.395,0)$ for the two-image, four-image,
 near-fold, and near-cusp panels, respectively.  The curves are evaluated with
 the source-plane operator representation of
 Sec.~\ref{sec:operator-total}.}
 \label{fig:amplification}
\end{figure*}

At low frequency the Fresnel scale samples a broad part of the Fermat surface,
and the response changes slowly.  At larger $w$, localized stationary regions
interfere.  The fringe spacing is governed by image-delay differences, while
the contrast is governed by the relative magnifications.  Several delay
differences contribute inside the four-image region.

\subsection{Evaluation from the point-mass solution}
\label{sec:operator-total}

Equation~\eqref{eq:diffraction-integral} is the definition of the
amplification, but the calculations in this paper do not evaluate it by
lens-plane quadrature.  The external shear is quadratic in the lens-plane
coordinates and can be represented by derivatives with respect to the source
position.

First set $\kappa=0$ and remove the source phase,
\begin{equation}
 F_{0,\gamma}(w,\by;\beta)
 =e^{\ii w|\by|^2/2}\cG_{\gamma}(w,\by;\beta).
 \label{eq:G-definition}
\end{equation}
The exact operator identity is
\begin{equation}
 \cG_{\gamma}(w,\by;\beta)
 =\exp\!\left[\frac{\ii\gamma}{2w}D_\beta\right]
  \cG_{\rm PM}(w,\by),
 \label{eq:operator-identity}
\end{equation}
where
\begin{equation}
 D_\beta=
 \cos 2\beta\,(\partial_1^2-\partial_2^2)
 +2\sin 2\beta\,\partial_1\partial_2.
 \label{eq:D-beta}
\end{equation}
Writing $s=|\by|^2$, the point-mass seed is
\begin{align}
 \cG_{\rm PM}(w,s)
 ={}&C(w)\,{}_1F_1\!\left(1-\frac{\ii w}{2};1;
                  -\frac{\ii ws}{2}\right),
 \label{eq:point-mass-seed}\\
 C(w)={}&
 \exp\!\left[\frac{\pi w}{4}
 +\frac{\ii w}{2}\ln\!\left(\frac{w}{2}\right)\right]
 \Gamma\!\left(1-\frac{\ii w}{2}\right).
 \label{eq:point-mass-prefactor}
\end{align}
Every source derivative generated by $D_\beta$ is another confluent
hypergeometric function.  The operator exponential can therefore be evaluated
as an adaptively truncated series.  Appendix~\ref{app:operator} gives the
derivation and the radial-derivative formula used in the calculation.

External convergence is included by an exact rescaling.  Define
\begin{equation}
 \lambda=1-\kappa,
 \qquad
 \by'=\frac{\by}{\sqrt{\lambda}},
 \qquad
 \gamma'=\frac{\gamma}{\lambda}.
 \label{eq:mass-sheet-vars}
\end{equation}
Changing variables to $\bx'=\sqrt{\lambda}\,\bx$ in
Eq.~\eqref{eq:diffraction-integral} gives
\begin{equation}
 F_{\kappa,\gamma}(w,\by;\beta)
 =\frac{1}{\lambda}
 \exp\!\left[
  \frac{\ii w}{2}\ln\lambda
  -\frac{\ii w\kappa}{2\lambda}|\by|^2
 \right]
 F_{0,\gamma'}(w,\by';\beta).
 \label{eq:mass-sheet-map}
\end{equation}
The dependence on $\kappa$ is therefore algebraic once the pure-shear series
has been evaluated.  No numerical contour integral appears in the final
procedure.

\section{Stationary phase and the origin of the fringes}
\label{sec:stationary}

The image positions satisfy
\begin{equation}
 \nabla_{\bx}\tau=0,
 \qquad
 \by=\bA\bx-\frac{\bx}{|\bx|^2}.
 \label{eq:lens-equation}
\end{equation}
Let $\bx_j$ be a nondegenerate solution, $\tau_j=\tau(\bx_j;\by)$, and
\begin{equation}
 \bH_j=\left.\nabla_{\bx}\nabla_{\bx}\tau\right|_{\bx_j}
 \label{eq:fermat-hessian}
\end{equation}
its Hessian.  The signed magnification is $\mu_j=(\det\bH_j)^{-1}$, and the
Morse index $n_j$ is the number of negative eigenvalues of $\bH_j$.  Applying
stationary phase gives
\begin{equation}
 F(w,\by)\sim\sum_j e^{\ii w\tau_j}H_j(w),
 \label{eq:image-sum}
\end{equation}
with
\begin{equation}
 H_j(w)=\sqrt{|\mu_j|}\,e^{-\ii\pi n_j/2}
 \left[1+\frac{\ii C_{1j}}{w}
       +\frac{C_{2j}}{w^2}+O(w^{-3})\right].
 \label{eq:image-kernel}
\end{equation}
The derivation and the explicit coefficients $C_{1j}$ and $C_{2j}$ are given
in Appendix~\ref{app:stationary-coefficients}.  The higher-order terms improve
the resolved-image approximation after the asymptotic ordering has set in.

Figure~\ref{fig:stationary-phase} compares the operator result with the leading
image sum and the expansion through $w^{-2}$ in regular two-image and
four-image regions.  The four-image source was chosen so that none of the
images is close to a fold.  A separate example in
Sec.~\ref{sec:localchannels} shows that the isolated-image expansion can remain
poor at moderate frequency even when the source is formally inside a
four-image region.

\begin{figure*}[t]
 \centering
 \includegraphics[width=0.94\textwidth]{figures/stationary_phase_comparison.pdf}
 \caption{Stationary-phase approximation in regular two-image and four-image
 regions for $\kappa=\gamma=0.2$ and $\beta=0$.  The source positions are
 $\by=(0.90,0)$ and $(0.03,0.10)$, respectively.  The plotted interval begins at
 $w=12$, where the isolated-image expansion is asymptotically ordered for the
 selected sources.  The $w^{-2}$ result closely follows the operator total.}
 \label{fig:stationary-phase}
\end{figure*}

For two images, Eq.~\eqref{eq:image-sum} gives
\begin{align}
 |F|^2\simeq{}&|H_1|^2+|H_2|^2
 +2\,\mathrm{Re}\!\left[
 H_1H_2^*e^{\ii w(\tau_1-\tau_2)}\right].
 \label{eq:two-image-interference}
\end{align}
If the kernels vary slowly over a fringe, the period is
\begin{equation}
 \Delta w\simeq\frac{2\pi}{|\tau_1-\tau_2|}.
 \label{eq:fringe-period}
\end{equation}
The same interpretation applies to every image pair in a multi-image region.
The delay phases carry the rapid frequency dependence that must be separated
before interpolation.

\section{Why ordinary relative binning fails}
\label{sec:rbfailure}

For stationary Gaussian detector noise, define the complex overlap
\begin{equation}
 \cZ[a,b]=4\sum_f\frac{a(f)b^*(f)}{S_n(f)}\,\Delta f.
 \label{eq:complex-overlap}
\end{equation}
The log likelihood is
\begin{equation}
 \ln\mathcal L=\mathrm{Re}\,\cZ[d,h]
 -\frac12\cZ[h,h]+\mathrm{const}.
 \label{eq:likelihood}
\end{equation}
Relative binning chooses a fiducial waveform $h_0$ and writes
\begin{equation}
 h(f;\vartheta)=r(f;\vartheta)h_0(f).
 \label{eq:rb-ratio}
\end{equation}
When $r$ is smooth, it can be represented by a piecewise low-order polynomial,
while the sums involving $d$ and $h_0$ are precomputed
\cite{ZackayDaiVenumadhav2018}.

For a lensed waveform,
\begin{equation}
 r(f)=
 \frac{\widetilde h_{\rm U}(f;\lambda)}
      {\widetilde h_{\rm U}(f;\lambda_0)}
 \frac{F[w(f);\vartheta_{\rm L}]}
      {F[w_0(f);\vartheta_{{\rm L},0}]},
 \label{eq:lensed-ratio}
\end{equation}
where $\lambda$ denotes binary parameters and $\vartheta_{\rm L}$ denotes lens
parameters.  The second factor need not be smooth.  For two images,
\begin{equation}
 F=e^{\ii w\tau_1}
 \left[H_1+H_2e^{\ii w\Delta\tau}\right],
 \qquad
 \Delta\tau=\tau_2-\tau_1,
 \label{eq:two-image-factorization}
\end{equation}
so the lens ratio contains interference patterns with different delays in the
candidate and fiducial models.  Near a minimum of the fiducial amplification,
the ratio can also acquire a sharp amplitude peak and a rapid phase excursion.

The loss of smoothness therefore occurs within a region of fixed image
topology.  It follows from the coherent addition of contributions with
different parameter-dependent delays, and does not require the creation or
destruction of images.  A numerical comparison of this direct ratio with the
separated representation is given after the latter has been constructed in
Sec.~\ref{sec:numerics}.

\section{Projected image contributions at fixed topology}
\label{sec:localchannels}

The stationary-phase form suggests retaining the image delays as explicit
phases.  Within a region where the images can be labeled continuously, write
\begin{equation}
 F(w;\vartheta_{\rm L})
 =\sum_{a=1}^{N_{\rm img}}
 e^{\ii w\tau_a(\vartheta_{\rm L})}
 K_a(w;\vartheta_{\rm L}).
 \label{eq:local-channel-decomposition}
\end{equation}
The functions $K_a$ will be called image contributions in this section.  They
approach the stationary-phase kernels $H_a$ when the images are well resolved.

At finite frequency, using $K_a=H_a$ would impose the error of the
stationary-phase expansion on the total.  Instead, let $\widehat K_a$ be any
set of trial contributions and define
\begin{equation}
 R(w)=F_{\rm op}(w)-\sum_b e^{\ii w\tau_b}\widehat K_b(w),
 \label{eq:local-residual}
\end{equation}
where $F_{\rm op}$ is the converged operator result.  Choose smooth weights
$\alpha_a$ with $\sum_a\alpha_a=1$ and set
\begin{equation}
 \boxed{
 K_a(w)=\widehat K_a(w)+
 \alpha_a e^{-\ii w\tau_a}R(w).}
 \label{eq:local-projection}
\end{equation}
Substitution into Eq.~\eqref{eq:local-channel-decomposition} gives
$F=F_{\rm op}$ at every frequency.  In a resolved region one may take
$\widehat K_a=H_a$.  The residual then contains the finite-frequency part that
is not described by the isolated saddles.  Equation~\eqref{eq:local-projection}
is useful even when the image topology never changes.

Figure~\ref{fig:projected-residual} uses a four-image source for which the
isolated-image expansion is poorly ordered at moderate $w$.  The projected sum
agrees with the operator total to floating-point accuracy.  The stationary-
phase terms still determine the physical high-frequency limit, while the
projection supplies the missing finite-frequency contribution.

\begin{figure*}[t]
 \centering
 \includegraphics[width=0.91\textwidth]{figures/projected_residual_comparison.pdf}
 \caption{A four-image example at $\by=(0.10,0.10)$,
 $\kappa=\gamma=0.2$.  The isolated-image expansion through $w^{-2}$ is
 inaccurate at moderate frequency because two of the saddles have small
 Hessian eigenvalues.  The projected image sum uses the same saddle terms as
 targets and adds the residual in Eq.~\eqref{eq:local-projection}; its
 reconstruction error is below $9\times10^{-16}$ in this calculation.}
 \label{fig:projected-residual}
\end{figure*}

Within a region of fixed topology, the candidate and fiducial waveforms can
therefore be compared image by image.  The difference between their image
delays is retained analytically, while only the ratio of the projected
amplitudes is interpolated.  This construction addresses the interference
fringes without relying on the accuracy of the isolated-image expansion at
moderate frequency.  Its remaining limitation is geometric: a physical image
label cannot be followed through a caustic at which that image is created or
destroyed.

\section{Changes in image multiplicity}
\label{sec:topology}

For the positive-parity Chang--Refsdal lens, a source outside the astroid
caustic has two nondegenerate images and a source inside has four
\cite{AnEvans2006}.  Crossing a fold creates or destroys a pair.  At a cusp,
three images meet at a single critical point.  The Hessian determinant of the
participating images vanishes at the caustic, so the individual saddle kernels
diverge even though the wave-optics amplification remains finite.

The left panel of Fig.~\ref{fig:topology-geometry} shows the fold and cusp
trajectories used in the numerical tests.  The signed path coordinate $\eta$
is negative outside the caustic and positive inside for the conventions used
here.

A fixed representation requires labels that exist before the real images do.
For each unoccupied label, we use the nearest point on the critical curve as a
geometric marker and assign the delay at that point.  When images are created,
their positions emerge from the marker and inherit its labels.  The delay
assigned to a label is kept at the critical-point delay while the cluster is
unresolved.  It approaches the physical image delay only after the delay
separation can be resolved within the analysis band.

For a cluster with physical delays $\tau_j^{\rm img}$ and critical delay
$\tau_c$, define the delay spread
$\Delta\tau_{\rm cl}=\max_j\tau_j^{\rm img}-\min_j\tau_j^{\rm img}$ and
\begin{equation}
 B_c=S_c(w_{\max}\Delta\tau_{\rm cl}),
 \qquad
 \tau_j^{\rm tr}=(1-B_c)\tau_c+B_c\tau_j^{\rm img}.
 \label{eq:tracked-delays}
\end{equation}
Here $S_c$ is a smooth transition function.  Outside the caustic all labels in
the cluster have $\tau_j^{\rm tr}=\tau_c$.  The right panels of
Fig.~\ref{fig:topology-geometry} show the resulting delays for the selected
paths.  The tracked delays remain coincident through the unresolved
neighborhood and then separate continuously toward the physical delays.

\begin{figure*}[t]
 \centering
 \begin{minipage}[c]{0.30\textwidth}
  \centering
  \includegraphics[width=\linewidth]{figures/caustic_paths.pdf}
 \end{minipage}\hfill
 \begin{minipage}[c]{0.66\textwidth}
  \centering
  \includegraphics[width=\linewidth]{figures/topology_stable_delays.pdf}
 \end{minipage}
 \caption{Source trajectories and delay continuation across the
 Chang--Refsdal caustic.  The left panel shows the astroid caustic at
 $\gamma=0.2$ and the fold and cusp paths used in the numerical tests.  The
 right panels show the physical image delays and the continuously tracked
 delays along those paths, measured relative to the critical-point delay.  The
 vertical lines mark the caustic.  The tracked delays remain coincident while
 the cluster is unresolved and approach the physical image delays farther
 inside the caustic.}
 \label{fig:topology-geometry}
\end{figure*}
\FloatBarrier

\section{Topology-stable four-component representation}
\label{sec:globalchannels}

The Chang--Refsdal lens has at most four nondegenerate images in the domain
considered here.  We retain four labels at every parameter point and continue
them using their lens-plane positions.  Labels without real images are placed
at the critical point associated with the caustic cluster from which those
images are created.  The continuation procedure is described in
Appendix~\ref{app:tracking}.

The exact projection is most useful when it is applied to the image cluster
that is poorly described by isolated stationary points.  Let $C$ denote the
active cluster and let $P$ denote the remaining, persistent images.  A generic
fold has two labels in $C$, while a cusp has three.  The persistent images are
assigned their stationary-phase factors $H_p$.  Their coherent sum is removed
from the operator result, leaving the exact contribution associated with the
cluster,
\begin{equation}
 G_C(w)=F_{\rm op}(w)-\sum_{p\in P}e^{\ii w\tau_p}H_p(w).
 \label{eq:exact-cluster-total}
\end{equation}
Any finite-frequency error in the persistent-image expansion is included in
$G_C$ by construction.  This choice is preferable to distributing the same
residual uniformly among all four labels, since it leaves the well-resolved
image factors close to their physical stationary-phase limits.

For each cluster label $j\in C$, choose smooth weights $\alpha_j$ satisfying
$\sum_{j\in C}\alpha_j=1$.  Equal weights are used in the calculations below.
When the cluster is unresolved, define
\begin{equation}
 L_j(w)=\alpha_j e^{-\ii w\tau_j}G_C(w).
 \label{eq:unresolved-cluster-split}
\end{equation}
The coherent sum of the $L_j$ is exactly $G_C$ for any assigned delays.  If a
cluster label corresponds to a real image, let $H_j$ be its stationary-phase
factor; for a virtual label set $H_j=0$.  For a real cluster image define
\begin{equation}
 \delta_j=\min_{k\in C,\,k\ne j}|\tau_j-\tau_k|
 \label{eq:delay-separation}
\end{equation}
and use the smooth switch
\begin{align}
 S_j(w)&=
 \begin{cases}
 0, & w\delta_j\leq\rho_0,\\[1mm]
 10u_j^3-15u_j^4+6u_j^5,
 & \rho_0<w\delta_j<\rho_1,\\[1mm]
 1, & w\delta_j\geq\rho_1,
 \end{cases}
 \label{eq:switch}\\
 u_j&=\frac{w\delta_j-\rho_0}{\rho_1-\rho_0}.
 \label{eq:switch-coordinate}
\end{align}
For a virtual label, $S_j=0$.  We use
$(\rho_0,\rho_1)=(0.5,4)$ in the calculations below.  The trial cluster
factors are
\begin{equation}
 \widehat K_j=(1-S_j)L_j+S_jH_j.
 \label{eq:trial-cluster}
\end{equation}
Their coherent sum need not equal $G_C$.  The discrepancy is projected back
onto the same cluster labels,
\begin{equation}
 \boxed{
 K_j=\widehat K_j+\alpha_j e^{-\ii w\tau_j}
 \left[G_C-\sum_{k\in C}e^{\ii w\tau_k}\widehat K_k\right],
 \qquad j\in C .}
 \label{eq:global-projection}
\end{equation}
For the persistent labels we set
\begin{equation}
 K_p=H_p,\qquad p\in P.
 \label{eq:persistent-factors}
\end{equation}
Equations~\eqref{eq:exact-cluster-total}--\eqref{eq:persistent-factors}
therefore give
\begin{equation}
 F_{\rm op}(w)=\sum_{a=1}^{4}e^{\ii w\tau_a}K_a(w)
 \label{eq:exact-reconstruction}
\end{equation}
at every frequency.

The individual cluster factors can be written without the intermediate trial
terms.  Define
\begin{equation}
 \overline S_C=\sum_{k\in C}\alpha_kS_k,
 \qquad
 Q_C=\sum_{k\in C}S_k e^{\ii w\tau_k}H_k.
 \label{eq:Sbar-Q}
\end{equation}
Then
\begin{equation}
 \boxed{\begin{aligned}
 K_j={}&S_jH_j+\alpha_j e^{-\ii w\tau_j}\\
 &\times\left[(1-S_j+\overline S_C)G_C-Q_C\right],
 \qquad j\in C .
 \end{aligned}}
 \label{eq:explicit-K}
\end{equation}
When the cluster is unresolved, $S_j=0$ and the exact residual is divided
among the coincident cluster labels.  Once the images are resolved and the
stationary-phase expansion becomes accurate, the correction in
Eq.~\eqref{eq:global-projection} decreases and $K_j$ approaches $H_j$.

The labels and the active cluster are continued along a path in lens-parameter
space.  At a fold, the two cluster labels remain present on both sides of the
caustic, although they correspond to real images only on the four-image side.
Near a cusp, three labels are associated with the coalescing triplet.  The
construction retains a fixed total of four components while localizing the
finite-frequency residual in the part of the image configuration for which
the isolated-saddle description is least reliable.

External convergence requires no change to this projection.  The operator
total is obtained from Eq.~\eqref{eq:mass-sheet-map}; the delays, Hessians, and
image positions are computed with the full matrix $\bA$.

\section{Multi-component relative binning}
\label{sec:multichannel-rb}

The global decomposition turns the lensed waveform into a sum of four
fiducial components.  Define
\begin{equation}
 h_{a0}(f)=\widetilde h_{\rm U}(f;\lambda_0)
 e^{\ii w_0(f)\tau_{a0}}
 K_{a0}[w_0(f)].
 \label{eq:fiducial-component}
\end{equation}
For a candidate lens and binary waveform,
\begin{equation}
 \widetilde h_{\rm L}(f)=\sum_a r_a(f)h_{a0}(f),
 \label{eq:rb-component-sum}
\end{equation}
where
\begin{align}
 r_a(f)={}&
 \frac{\widetilde h_{\rm U}(f;\lambda)}
      {\widetilde h_{\rm U}(f;\lambda_0)}
 \exp\!\left\{\ii\left[w(f)\tau_a-w_0(f)\tau_{a0}\right]\right\}
 \nonumber\\
 &\times\frac{K_a[w(f)]}{K_{a0}[w_0(f)]}.
 \label{eq:component-ratio}
\end{align}
This expression allows the point-mass redshifted mass to vary, in which case
$w(f)$ differs between the candidate and fiducial models.  The projected
residual in Eq.~\eqref{eq:global-projection} is already contained in $K_a$.
There is no separate residual term in the likelihood.

Divide the frequency range into bins centered at $f_b$, and approximate each
component ratio within a bin by
\begin{equation}
 r_a(f)\simeq r_{a,b}^{(0)}+r_{a,b}^{(1)}(f-f_b).
 \label{eq:linear-ratio}
\end{equation}
The data-component summaries are
\begin{equation}
 A_{a,b}^{(p)}=
 4\sum_{f\in b}\frac{d(f)h_{a0}^*(f)}{S_n(f)}(f-f_b)^p\Delta f,
 \qquad p=0,1.
 \label{eq:A-summary}
\end{equation}
They give
\begin{equation}
 \cZ[d,h]\simeq
 \sum_{a,b}\left[A_{a,b}^{(0)}r_{a,b}^{(0)*}
                  +A_{a,b}^{(1)}r_{a,b}^{(1)*}\right].
 \label{eq:data-term-rb}
\end{equation}

The waveform norm contains pairs of components.  Define
\begin{equation}
 B_{ac,b}^{(p)}=
 4\sum_{f\in b}\frac{h_{a0}(f)h_{c0}^*(f)}{S_n(f)}
 (f-f_b)^p\Delta f,
 \qquad p=0,1,2.
 \label{eq:B-summary}
\end{equation}
Multiplying the two linear expansions gives
\begin{align}
 \cZ[h,h]\simeq\sum_{a,c,b}\bigg\{&
 B_{ac,b}^{(0)}r_{a,b}^{(0)}r_{c,b}^{(0)*}\nonumber\\
 &+B_{ac,b}^{(1)}\left[
 r_{a,b}^{(1)}r_{c,b}^{(0)*}
 +r_{a,b}^{(0)}r_{c,b}^{(1)*}\right]\nonumber\\
 &+B_{ac,b}^{(2)}r_{a,b}^{(1)}r_{c,b}^{(1)*}\bigg\}.
 \label{eq:norm-term-rb}
\end{align}
The Hermitian relation $B_{ca,b}^{(p)}=B_{ac,b}^{(p)*}$ reduces the number of
independent component pairs.  For four labels, ten pair summaries are required
per moment.  The factor-of-several increase in precomputed summaries is small
compared with the reduction in frequency nodes when the complete amplification
ratio is strongly oscillatory.

Equation~\eqref{eq:component-ratio} separates two kinds of frequency
dependence.  The ratio $K_a/K_{a0}$ contains the finite-frequency response,
including the projected residual.  The exponential contains the known
difference between the candidate and fiducial image delays.  We retain this
exponential analytically and interpolate only the unlensed-waveform ratio and
$K_a/K_{a0}$.  Since $w$ is proportional to frequency, write
\begin{equation}
 w(f)\tau_a-w_0(f)\tau_{a0}=2\pi f\,\delta t_a
 \label{eq:channel-time-shift}
\end{equation}
with a frequency-independent relative delay $\delta t_a$, and define
\begin{equation}
 q_a(f)=
 \frac{\widetilde h_{\rm U}(f;\lambda)}
      {\widetilde h_{\rm U}(f;\lambda_0)}
 \frac{K_a[w(f)]}{K_{a0}[w_0(f)]}.
 \label{eq:slow-component-ratio}
\end{equation}
Then $r_a=e^{2\pi\ii f\delta t_a}q_a$.  Approximating only $q_a$ within each
bin replaces the summaries in Eqs.~\eqref{eq:A-summary} and
\eqref{eq:B-summary} by the same sums evaluated at the time shifts
$\delta t_a$ and $\delta t_a-\delta t_c$, respectively.  These shifted sums
are short discrete Fourier transforms and can be tabulated on a delay grid or
evaluated with FFTs.  This phase-aware form is useful when the candidate and
fiducial delays differ enough for the exponential to accumulate several cycles
across the detector band.  The numerical benchmarks below retain this delay phase exactly and
interpolate only $q_a$.

\section{Numerical results}
\label{sec:numerics}

The tests use $\gamma=0.2$ and $5\leq w\leq40$ unless stated otherwise.  The
operator series is truncated after four successive terms fall below
$2\times10^{-12}$ times the partial sum, with a maximum order chosen high
enough that every evaluation in the reported local tests satisfies this
criterion.  The image geometry is obtained from the quartic reduction in
Appendix~\ref{app:quartic}.  The numerical interpolation diagnostic is
\begin{equation}
 \epsilon(w)=\frac{|F_{\rm int}(w)-F_{\rm op}(w)|}
 {\max\left[|F_{\rm op}(w)|,
 0.15\max_{w'}|F_{\rm op}(w')|\right]}.
 \label{eq:error-profile}
\end{equation}
The floor prevents a zero of the exact amplification from dominating the
reported maximum.  It does not enter the channel construction.

\subsection{Fixed-topology comparison}

Figure~\ref{fig:projected-ratios} compares the two representations for
sources that remain outside the same generic fold.  Both models have two real
images, yet the ratio of their complete amplification factors develops a
sequence of narrow peaks and phase excursions because the interference minima
occur at different frequencies.  Once the four contributions are separated,
the ratios of their complex amplitudes vary slowly across the same interval.  The known delay differences are treated
analytically and are therefore absent from the lower panels.

\begin{figure*}[t]
 \centering
 \includegraphics[width=0.92\textwidth]{figures/projected_kernel_ratios.pdf}
 \caption{Direct and separated candidate-to-fiducial ratios within a fixed
 two-image topology.  The two sources lie outside the same generic fold at
 signed normal displacements $\eta_s=-0.040$ and $\eta_s=-0.010$, with
 $\gamma=0.2$.  The complete amplification ratio in the upper panels develops a sequence of
 sharp peaks and phase excursions as the interference minima shift.  The
 lower panels show the corresponding ratios of the two persistent image
 amplitudes and the two virtual fold-cluster amplitudes.  Their phases are
 measured relative to the lowest plotted frequency.  The known differences
 between the candidate and fiducial delays are omitted from the lower panels
 and retained analytically when the four contributions are recombined.}
 \label{fig:projected-ratios}
\end{figure*}

A greedy interpolation test starts with the two endpoints and repeatedly adds
the frequency with the largest value of Eq.~\eqref{eq:error-profile}.  The
differences between the candidate and fiducial delays are retained exactly;
only the slowly varying factor ratios shown in the lower panels of
Fig.~\ref{fig:projected-ratios} are interpolated.  The separated
representation reaches $\max_w\epsilon<10^{-3}$ with 11 frequency nodes.
Direct interpolation of $F/F_0$ requires 135 of the 141 validation
frequencies for the same criterion.

\subsection{Fold and cusp crossings}

For the local crossing tests, two persistent contributions and a two-term
cluster are used at a fold; one persistent contribution and a three-term
cluster are used at a cusp.  These use the cluster-local projection in
Eqs.~\eqref{eq:exact-cluster-total}--\eqref{eq:global-projection}.  Figure~\ref{fig:crossing-ratios}
shows the candidate-to-fiducial ratios for complete outside-to-inside
crossings.  Several cluster curves coincide because equal weights are used while the
corresponding delays are unresolved.

\begin{figure*}[t]
 \centering
 \includegraphics[width=0.92\textwidth]{figures/channel_ratios.pdf}
 \caption{Amplitude and unwrapped phase of the separated ratios for a fold
 crossing and a cusp crossing from $\eta=-0.01$ to $\eta=+0.01$.  The ratios
 remain smooth even though the number of real images changes.}
 \label{fig:crossing-ratios}
\end{figure*}

Table~\ref{tab:nodes} gives the number of adaptively selected frequency nodes
needed to reach $\max_w\epsilon<10^{-3}$.  The separated representation needs
7--10 nodes for the four representative crossings.  Direct interpolation of
the complete ratio requires 70 or 71 of the 71 validation frequencies.
Across all crossing pairs in the scan $-0.02\leq\eta\leq0.02$, the separated
representation uses 6--10 nodes.

\begin{table}[t]
\caption{Frequency nodes required for a maximum normalized interpolation error
of $10^{-3}$ over $5\leq w\leq40$.}
\label{tab:nodes}
\centering
\begin{tabular}{lcc}
\toprule
Crossing & separated ratios & complete ratio\\
\midrule
Fold, outside to inside & 9 & 70\\
Fold, inside to outside & 10 & 70\\
Cusp, outside to inside & 7 & 71\\
Cusp, inside to outside & 7 & 71\\
\bottomrule
\end{tabular}
\end{table}

Figure~\ref{fig:adaptive-nodes} shows the error as the node budget is
increased.  The direct-ratio curves remain above the target on the displayed
node scale because they must resolve the moving fringes.

\begin{figure}[t]
 \centering
 \includegraphics[width=0.98\columnwidth]{figures/adaptive_nodes.pdf}
 \caption{Maximum interpolation error versus frequency-node budget for a full
 fold crossing and a full cusp crossing.  The dashed line marks the
 $10^{-3}$ target.}
 \label{fig:adaptive-nodes}
\end{figure}

The coherent component sum reproduces the operator total to between
$4.4\times10^{-16}$ and $6.7\times10^{-16}$ in the four local crossing
calculations.  The smallest component magnitudes remain above 0.47 in the cusp
tests and 0.59 in the fold tests, so the candidate-to-fiducial ratios are not
made ill conditioned by a vanishing denominator.

\subsection{Geometry, convergence, and global tracking}

The quartic image solver was compared with a deterministic two-dimensional
multistart solver on 168 configurations, including general source-shear
orientations and points within $10^{-5}$ of selected folds and cusps.  The
image count agreed in every case.  The maximum matched position difference was
$2.1\times10^{-12}$, and the maximum lens-equation residual of a quartic image
was $1.8\times10^{-13}$.  In this scan the quartic calculation was about 190
times faster than the multistart calculation.  The timing indicates that image
finding is negligible compared with the special-function evaluation in the
current prototype.

The external-convergence mapping was tested for six values between
$\kappa=0.2$ and $0.3$, while keeping the reduced shear and scaled source
position fixed.  The four-term reconstruction errors ranged from
$2.0\times10^{-15}$ to $6.3\times10^{-15}$.

A separate continuation path varied the source position, shear magnitude, and
shear orientation through two-image and four-image regions.  The path contains
four changes in image count.  All four labels remained defined.  The largest
absolute reconstruction error over the 49 path points was
$9.4\times10^{-11}$; this occurred where large intermediate terms produced
strong floating-point cancellation.  Away from those points the errors were
near $10^{-15}$.  The continuation algorithm is designed for nearby parameter
points.  A distant independent proposal should be initialized from a fixed
convention or connected to the fiducial point by a short path.

\FloatBarrier
\section{Discussion}
\label{sec:discussion}

The failure of ordinary relative binning in wave-optics lensing is caused by
the coexistence of several parameter-dependent image phases.  It does not
require a caustic crossing.  Separating those phases produces functions whose
candidate-to-fiducial ratios are much smoother.  The residual projection in
Eq.~\eqref{eq:local-projection} makes this statement useful at finite
frequency, where the isolated-image expansion may be inaccurate.  The
stationary-phase approximation supplies the physical resolved-image limit,
while the operator result fixes the coherent sum at every frequency.

The topology-stable extension addresses a distinct issue.  Physical images
cannot be labeled globally because pairs are created at folds and triplets
meet at cusps.  Four continuously tracked labels are sufficient for a
positive-parity Chang--Refsdal lens.  Virtual labels are placed at a critical
point, and the representation remains finite while their physical images are
absent.  Once the images are resolved, the corresponding terms approach the
usual geometrical contributions.

The numerical tests are intended to establish the representation and its
interpolation behavior over a controlled domain.  They do not yet constitute
a full parameter-estimation study.  The source positions, convergence, and
shear were varied, while the interpolation tables in
Table~\ref{tab:nodes} keep the redshifted microlens mass fixed.  Equation
\eqref{eq:component-ratio} includes mass variation and ordinary binary
parameters, so these can be incorporated in the same relative-binning
summaries.  Their joint effect on bin placement should be tested with detector
noise curves and complete compact-binary waveform models.

The present operator implementation is restricted to positive-parity macro
images with $1-\kappa>|\gamma|$.  Macro saddles are important for strongly
lensed events and require a separate treatment.  The source-derivative series
also becomes more demanding as the reduced shear approaches unity or as the
frequency is extended far beyond the range studied here.  Stable recurrences
or a reduced representation of the operator total would be useful in that
regime.

The decomposition is not specific to relative binning.  It provides a set of
frequency functions with explicit rapidly varying phases and smooth parameter
dependence.  The same representation may be useful for surrogate models,
reduced-order quadrature, and interpolation of microlensing amplification
factors.  For the application considered here, its immediate consequence is a
multi-component relative-binning likelihood that remains well defined across
the Chang--Refsdal caustic network and contains no numerical lens-plane
integral.

\section{Conclusion}
\label{sec:conclusion}

We have constructed a decomposition of the positive-parity Chang--Refsdal
amplification factor in which the image-delay phases are retained explicitly.
Within a region of fixed topology, an exact residual projection supplements
the isolated-image stationary-phase terms and reproduces the complete
wave-optics result.  Four continuously tracked labels extend the decomposition
through folds and cusps.  External convergence is included by an exact
mass-sheet rescaling, and the all-frequency amplification is evaluated from
the point-mass solution through a source-derivative expansion.

The separated candidate-to-fiducial ratios are smooth in the fixed-topology
and caustic-crossing examples studied here.  Complete fold and cusp crossings
reach a maximum interpolation error of $10^{-3}$ with 6--10 frequency nodes,
whereas direct interpolation of the total amplification ratio requires nearly
the full validation grid.  The likelihood formulas in
Eqs.~\eqref{eq:data-term-rb} and~\eqref{eq:norm-term-rb} show how these ratios
enter relative binning without a separate treatment of the projected
finite-frequency residual.

\section*{Data and code availability}

The scripts, numerical data, and contour-free implementation used to produce
the figures and tables accompany this manuscript.  The numerical lens-plane
integral is not used by the code path described in the paper.

\appendix

\section{Quartic image geometry}
\label{app:quartic}

The image equation can be reduced to a quartic polynomial.  This appendix
gives the reduction for a general symmetric matrix $\bA$, so convergence and
an arbitrarily oriented shear are included.

Let $Y=|\by|$ and rotate to the source-aligned basis
$\bm B=(\widehat\by,\widehat\by_\perp)$.  In this basis,
\begin{equation}
 \by'=(Y,0),
 \qquad
 \bA'=\bm B^{\mathsf T}\bA\bm B
 =\begin{pmatrix}a&b\\b&d\end{pmatrix}.
 \label{eq:A-prime}
\end{equation}
Set $u=1/|\bx'|^2$.  The lens equation becomes
\begin{equation}
 (\bA'-u\bm{1})\bx'=(Y,0).
 \label{eq:lens-u}
\end{equation}
With
\begin{equation}
 D(u)=(a-u)(d-u)-b^2,
\end{equation}
the generic solution for a fixed $u$ is
\begin{equation}
 x'_1=\frac{Y(d-u)}{D(u)},
 \qquad
 x'_2=-\frac{Yb}{D(u)}.
 \label{eq:quartic-reconstruction}
\end{equation}
Imposing $|\bx'|^2=u^{-1}$ gives
\begin{equation}
 D(u)^2-Y^2u\left[(d-u)^2+b^2\right]=0.
 \label{eq:quartic-compact}
\end{equation}
Expanding,
\begin{align}
0={}&u^4+[-2(a+d)-Y^2]u^3\nonumber\\
&+[a^2+4ad+d^2-2b^2+2dY^2]u^2\nonumber\\
&+[-2(a+d)(ad-b^2)-Y^2(d^2+b^2)]u
 +(ad-b^2)^2.
 \label{eq:quartic-expanded}
\end{align}
Positive real roots are reconstructed with
Eq.~\eqref{eq:quartic-reconstruction}, rotated back to the original frame, and
polished with deterministic Newton iterations.  The polishing corrects
floating-point error in the algebraic roots; it is not an image search.

When $b=0$, Eq.~\eqref{eq:quartic-reconstruction} has a removable singularity
for the symmetric off-axis pair.  The axial images solve
\begin{equation}
 a{x'_1}^2-Yx'_1-1=0,
\end{equation}
and an off-axis pair, when present, satisfies
\begin{equation}
 u=d,
 \qquad
 x'_1=\frac{Y}{a-d},
 \qquad
 {x'_2}^2=\frac1d-{x'_1}^2.
\end{equation}
At $Y=0$ and nonzero shear, the images lie along the eigenvectors of $\bA$ at
radii equal to the inverse square roots of its eigenvalues.  The exactly
unsheared and aligned case is an Einstein ring and is treated as a degeneracy.

\section{Stationary-phase coefficients}
\label{app:stationary-coefficients}

Write $\bx=\bx_j+\bm\xi$ near a nondegenerate image and expand
\begin{equation}
 \tau=\tau_j+\frac12H_{ab}\xi_a\xi_b
 +\frac1{3!}T_{abc}\xi_a\xi_b\xi_c
 +\frac1{4!}Q_{abcd}\xi_a\xi_b\xi_c\xi_d+\cdots.
 \label{eq:local-taylor}
\end{equation}
After the scaling $\bm\xi=w^{-1/2}\bm z$, the quadratic term is independent of
$w$ and the higher derivatives generate an expansion in powers of $w^{-1}$.
The oscillatory Gaussian moments are evaluated with covariance
$\ii\bH^{-1}$.  At first order, Wick contractions give
\begin{align}
 C_1={}&-\frac18 Q_{abcd}G_{ab}G_{cd}
 +\frac1{12}T_{abc}T_{def}G_{ad}G_{be}G_{cf}\nonumber\\
 &+\frac18T_{abc}T_{def}G_{ab}G_{de}G_{cf},
 \qquad G=\bH^{-1}.
 \label{eq:C1-contraction}
\end{align}
The second coefficient receives terms from derivatives through sixth order and
from products of the lower-order vertices.  For the Chang--Refsdal potential,
all third and higher derivatives arise from $-\ln|\bx|$.  This allows the
contractions to be reduced to a polynomial in three scaled inverse-Hessian
entries.

Let $\bm e_r=\bx_j/|\bx_j|$ and let $\bm e_t$ be the positively oriented
orthogonal vector.  Define
\begin{align}
 \bm M&=
 \begin{pmatrix}\bm e_r^{\mathsf T}\\\bm e_t^{\mathsf T}\end{pmatrix}
 \bH_j^{-1}
 \begin{pmatrix}\bm e_r&\bm e_t\end{pmatrix},
 \label{eq:M-def}\\
 u&=|\bx_j|^{-2},
 \qquad
 p=uM_{rr},\quad q=uM_{rt},\quad s=uM_{tt}.
 \label{eq:pqs}
\end{align}
Then
\begin{align}
C_1=\frac{1}{12}\big(&10p^3-12p^2s-9p^2-48pq^2+18ps^2+18ps\nonumber\\
&+72q^2s+36q^2-9s^2\big).
 \label{eq:C1-explicit}
\end{align}
The second coefficient is $C_2=-P_2/288$, with the polynomial below.

\begin{widetext}
\begin{align}
P_2={}&1540p^6-1680p^5s-3780p^5-16800p^4q^2
+2520p^4s^2+5040p^4s+2961p^4\nonumber\\
&+40320p^3q^2s+40320p^3q^2-3600p^3s^3-8280p^3s^2
-5364p^3s-720p^3\nonumber\\
&+40320p^2q^4-64800p^2q^2s^2-99360p^2q^2s-32184p^2q^2
+3780p^2s^4+10800p^2s^3\nonumber\\
&+9126p^2s^2+2160p^2s-86400pq^4s-66240pq^4
+60480pq^2s^3+129600pq^2s^2\nonumber\\
&+73008pq^2s+8640pq^2-3780ps^4-5940ps^3-2160ps^2
-11520q^6+60480q^4s^2\nonumber\\
&+86400q^4s+24336q^4-30240q^2s^3-35640q^2s^2
-8640q^2s+945s^4+720s^3.
 \label{eq:C2-poly}
\end{align}
\end{widetext}

Substitution of Eqs.~\eqref{eq:C1-explicit} and~\eqref{eq:C2-poly} into
Eq.~\eqref{eq:image-kernel} gives the resolved-image targets used in the main
text.

\section{Source-plane operator representation}
\label{app:operator}

For $\kappa=0$, write the shear contribution to the Fermat potential as
$-(\gamma/2)x_iQ_{ij}x_j$, where $Q$ is the traceless matrix in
Eq.~\eqref{eq:shear-matrix} divided by $\gamma$.  After removing the source
phase, the remaining lens-plane integral contains
$\exp(-\ii w\by\cdot\bx)$.  Acting on this Fourier kernel,
\begin{equation}
 \partial_{y_i}\partial_{y_j}e^{-\ii w\by\cdot\bx}
 =-w^2x_ix_j e^{-\ii w\by\cdot\bx}.
\end{equation}
Multiplication by the quadratic shear phase is therefore equivalent to the
source-space operator in Eq.~\eqref{eq:operator-identity}.  Expanding its
exponential gives
\begin{equation}
 \cG_{\gamma}=
 \sum_{n=0}^{\infty}\frac1{n!}
 \left(\frac{\ii\gamma}{2w}\right)^nD_\beta^n\cG_{\rm PM}.
 \label{eq:operator-series}
\end{equation}

The point-mass seed depends only on $s=|\by|^2$.  Its radial derivatives are
\begin{align}
\frac{\dd^k\cG_{\rm PM}}{\dd s^k}={}&
C(w)\left(-\frac{\ii w}{2}\right)^k
\frac{(1-\ii w/2)_k}{k!}\nonumber\\
&\times{}_1F_1\!\left(1-\frac{\ii w}{2}+k;1+k;
-\frac{\ii ws}{2}\right).
 \label{eq:radial-derivatives}
\end{align}
Repeated application of $D_\beta$ is represented as a finite sum of monomials
in $z=y_1+\ii y_2$, $\bar z$, and the derivatives in
Eq.~\eqref{eq:radial-derivatives}.  The series is stopped after a prescribed
number of consecutive terms fall below the requested fraction of the current
partial sum.  The order used and the last-term estimate are retained as
convergence diagnostics.

For completeness, the convergence rescaling follows by substituting
$\bx'=\sqrt{\lambda}\bx$ into Eq.~\eqref{eq:fermat}:
\begin{align}
 \tau_{\kappa,\gamma}(\bx;\by)={}&
 \tau_{0,\gamma/\lambda}
 \left(\bx';\frac{\by}{\sqrt\lambda}\right)
 +\frac12\ln\lambda
 -\frac{\kappa}{2\lambda}|\by|^2.
\end{align}
The Jacobian $\dd^2x=\lambda^{-1}\dd^2x'$ then gives
Eq.~\eqref{eq:mass-sheet-map}.

\section{Global label continuation}
\label{app:tracking}

At a given source position, the real images are found from
Eq.~\eqref{eq:quartic-expanded}.  The nearest caustic point is obtained by a
one-dimensional minimization along the critical curve.  Unoccupied labels are
placed at the corresponding critical point in the lens plane.

Suppose the previous parameter point has marker positions $\bm m_a$, one for
each of the four labels, and the current point has $N$ real images
$\bx_j$.  The assignment is chosen by minimizing
\begin{equation}
 \sum_{j=1}^{N}\frac{|\bx_j-\bm m_{\pi(j)}|^2}{\ell^2}
 \label{eq:assignment-cost}
\end{equation}
over injections $\pi$ from the images to the four labels.  The scale $\ell$ is
taken from the median marker radius and only makes the cost dimensionless.
Since there are at most four images, every assignment can be checked directly.
The markers of occupied labels are updated to the image positions; the others
are set to the nearest critical point.

At the first point on a path, the real images are ordered by polar angle and
then by radius.  This convention is sufficient to initialize a continuous
trajectory.  It does not define a globally path-independent permutation around
all possible loops in parameter space.  For likelihood proposals that are far
from the fiducial point, one may reset the labels with the initialization rule
or evaluate the model along a short interpolation path from the fiducial
parameters.  The coherent sum in Eq.~\eqref{eq:exact-reconstruction} is
independent of the label permutation; only the smoothness of individual ratios
depends on consistent continuation.

\bibliographystyle{apsrev4-2}
\bibliography{refs}

\end{document}
